(1) 热流从1到2为 $Q=G(T_1-T_2)$，故 $C_1\dot T_1=-G(T_1-T_2)$、$C_2\dot T_2=G(T_1-T_2)$。相加说明总显热 $C_1T_1+C_2T_2$ 恒定，温差满足 $\dot\Delta=-G(C_1^{-1}+C_2^{-1})\Delta$。

(2) 定义 $\bar T=(C_1T_{10}+C_2T_{20})/(C_1+C_2)$、$\tau=C_1C_2/[G(C_1+C_2)]$，则 $\Delta(t)=\Delta_0e^{-t/\tau}$。由能量平均及差分反求 $T_1(t)=\bar T+[C_2/(C_1+C_2)]\Delta(t)$、$T_2(t)=\bar T-[C_1/(C_1+C_2)]\Delta(t)$，最终同温为 $\bar T$。

(3) 总熵 $S=C_1\ln T_1+C_2\ln T_2+\text{常数}$，所以 $\dot S=C_1\dot T_1/T_1+C_2\dot T_2/T_2=G(T_1-T_2)^2/(T_1T_2)\ge0$。总熵增加量为 $C_1\ln(\bar T/T_{10})+C_2\ln(\bar T/T_{20})>0$，由加权算术—几何均值不等式得到。

(4) 可逆热机在最终共同温度 $T_{\rm rev}$ 时满足 $C_1\ln(T_{\rm rev}/T_{10})+C_2\ln(T_{\rm rev}/T_{20})=0$，故 $T_{\rm rev}=(T_{10}^{C_1}T_{20}^{C_2})^{1/(C_1+C_2)}\le\bar T$。最大可提取功为 $\boxed{W_{\max}=C_1T_{10}+C_2T_{20}-(C_1+C_2)T_{\rm rev}}$；直接导热因熵产使可提取功全部损失。
